% Appendix, from exam/README.md: what the two written papers are for, what they cover, the
% conventions they assume, and how they are marked.

\chapter{The Self-Assessment Papers}
\label{exam:ch:about}

This appendix and the four after it hold two complete three-hour papers for \emph{Zephyr RTOS},
with worked solutions. Appendices~\ref{exa:ch:paper} and \ref{exb:ch:paper} are the papers,
questions only; Appendices~\ref{ansa:ch:answers} and \ref{ansb:ch:answers} are their model answers,
with marks.

\section*{What these are for}
The book has \textbf{no exam and nothing is marked}. What you get instead is the exercise sets after
every chapter, each with a published solution, and the questions in every chapter's review.

These papers do not replace that. \textbf{They are here for you to test your own skills and
knowledge, nothing more.} They gate nothing, they are not a qualification, and no part of the book
requires them. Nothing in the companion repository depends on them, and neither \code{make examples}
nor \code{make test} knows they exist.

What they check is something the exercises cannot. Every exercise is done at a keyboard, with the
chapter open beside you and a build ready to say no. These papers ask what you can reconstruct with
none of that in front of you --- which of the rules you have absorbed, and which you have merely
been looking up.

\textbf{Take one after the last chapter.} Both papers draw on all six chapters, so sitting one
partway through examines material you have not read yet, and the result says more about how far you
have read than about what you have understood.

\section*{The two papers}
Both cover the whole book, Chapters~1 to 6, at the same weighting. They share no question. Either
can be used alone; use both as a main sitting and a resit, or in alternate years.

\begin{booktable}{@{}p{4.6em}Lp{5.4em}r@{}}
  \thead{Question} & \thead{Topic} & \thead{Chapters} & \thead{Marks} \\
  \midrule
  1 & The workspace and the build & 1 & 12 \\
  2 & Kconfig and logging & 1 & 12 \\
  3 & Devicetree, bindings and overlays & 2 & 13 \\
  4 & Threads, scheduling and time & 3 & 13 \\
  5 & Interrupt context, work queues and simulated time & 3 & 12 \\
  6 & Synchronization and communication & 4 & 13 \\
  7 & Writing a driver & 5 & 13 \\
  8 & Subsystems and debugging & 6 & 12 \\
  \midrule
    & & & \textbf{100} \\
\end{booktable}

\textbf{Both papers mix prose with artifacts}, because the subject does. Roughly half the marks are
for explaining what a build decided and why; the other half are for reading or writing the things a
Zephyr application is made of: a \code{prj.conf}, a devicetree overlay, a binding, a driver, a
thread.

The two papers differ in which half they lean on:
\begin{itemize}
  \item \textbf{Paper A leans towards what the build hands you.} Its questions put an artifact in
    front of you --- a build that aborted, a devicetree that produced the wrong macro, a trace of
    threads in the wrong order, a driver whose device is never ready --- and ask what happened and
    what to change. Naming the cause is half the marks; saying what it does to the running program
    is the other half.
  \item \textbf{Paper B leans towards what you write.} It asks for a Kconfig menu, a binding and the
    overlay that instantiates it, a sampling design split between an interrupt and a work queue, the
    skeleton of a devicetree-instantiated driver, and a ztest suite for it.
\end{itemize}

Neither paper asks anything that a careful reader could not answer from the chapters. Where a
question turns on a rule that is easy to state and easy to get backwards --- which file a Kconfig
value comes from last, why a log message that was filtered out is not in the binary, which context a
timer's expiry function runs in --- the solution says so at length, because that is where the marks
and the understanding both are.

\section*{Conventions the papers assume}
Both papers state these in their own rubric, so a candidate never has to have read this appendix.
\begin{itemize}
  \item \textbf{Zephyr 4.4 and the board \code{native_sim/native/64}}, built with the host's GCC, as
    in the book.
  \item \textbf{Code is marked on semantics, not syntax.} A missing \code{#include}, a missing
    semicolon or a brace in the wrong column costs nothing. A missing \code{device_is_ready()} check
    costs everything.
  \item \textbf{The book's house style is expected} where a question asks for it: Zephyr's own,
    meaning \code{snake_case}, the kernel's \code{k_} and \code{K_} spellings, application options
    prefixed \code{APP_}, \code{IS_ENABLED()} rather than \code{#ifdef} where both would work, and
    log modules rather than \code{printk()} in anything that is not a one-off.
  \item \textbf{The simulator's behaviour may be assumed} without deriving it: 100 ticks per second,
    a sleep that ends on the tick after the one it asked for, simulated time that only advances when
    every thread is waiting, immediate logging, and no thread analyzer, stack protection or
    userspace.
  \item \textbf{``State the consequence'' means the consequence.} Naming a cause earns half the
    marks; saying what it does to the running program earns the other half. ``This is wrong'' scores
    nothing.
  \item \textbf{Where a question asks for a stated number of defects}, listing more is not
    penalised, but only that many are marked, so the strongest answers go first.
\end{itemize}

\section*{Marking}
Every solution is written to be marked by somebody who has read the chapters and does not otherwise
write Zephyr applications for a living, so each carries the reasoning rather than the answer alone.
Marks are shown per part.

Three conventions worth agreeing before a paper is marked:
\begin{itemize}
  \item \textbf{Method carries the marks.} A correct diagnosis with a clumsy fix is worth more than
    a correct fix with no diagnosis. Several questions feed the next part, so an error should be
    followed through rather than penalised twice.
  \item \textbf{The named traps are worth full marks on their own.} A number of parts exist entirely
    to see whether a candidate walks into a specific mistake: a \code{CONFIG_} option misspelled in
    C rather than in \code{prj.conf}, a devicetree node that is \code{disabled} rather than absent, a
    \code{k_sleep()} in a timer's expiry function, a semaphore used as a lock, a driver that reads
    its configuration from \code{dev->data}, a busy loop on \code{native_sim}. Where a solution flags
    one of these, a candidate who walks into it loses those marks and no others.
  \item \textbf{The discussion parts are not decoration.} ``State why'' is where the book's actual
    content is, and a paper marked only on whether the code would build would pass a candidate who
    has understood none of it.
\end{itemize}
